\documentclass[10pt,twocolumn,letterpaper]{article}
\usepackage[T1]{fontenc}
\IfFileExists{cvpr.sty}{\usepackage[review]{cvpr}}{\usepackage[margin=0.8in]{geometry}}
\usepackage{amsmath,amssymb,booktabs,graphicx}
\usepackage[pagebackref,breaklinks,hidelinks]{hyperref}
\def\paperID{DRAFT}
\def\confName{CVPR}
\def\confYear{2027}
\title{VETO: Evidence-Constrained Harness Selection for Reliable Chart Reasoning}
\author{Anonymous authors}
\newcommand{\pending}{\textit{Pending experiment}}
\begin{document}
\maketitle
\begin{abstract}
Optimizing an executable harness for current task accuracy does not explicitly
test whether it responds correctly when answer-relevant visual evidence changes.
We study paired visual evidence as a selection criterion before a subsequent
weight update. VETO retains the incoming harness and selects among candidates
whose empirical paired correctness does not decrease. The weight update reuses
WHALE's native rejection-sampling fine-tuning. Our prospective study uses a
shared scoring protocol, matched candidate archives and independently seeded
training branches. \textbf{Study status: experiments are incomplete; this draft
does not claim a measured improvement.}
\end{abstract}

\section{Introduction}
Chart reasoning requires answers to follow the information in the presented
chart. Counterfactual chart evaluation already exposes failures hidden by
single-chart scores~\cite{chartographer}. Harness optimization and non-regression
acceptance also have established precedents~\cite{whale,autodesign}. We examine
the narrower question of whether paired evidence is useful for choosing the
harness that supplies subsequent successful training trajectories.

Our intended contribution is an experimentally supported selection mechanism,
not the invention of harness--weight alternation or counterfactual evaluation.
We use one complete training--search--training cycle and compare the paired gate
against a same-data single-image accuracy gate and counterfactual augmentation.
No long-horizon causal propagation or population reliability guarantee is assumed.

\section{Method}
\begin{figure*}[t]
\centering
\IfFileExists{method.pdf}{\includegraphics[width=.94\textwidth]{method.pdf}}{\fbox{Method figure pending build}}
\caption{Paired evidence constrains the harness selected for the next native RSFT stage.
The shared scorer, candidate search and weight optimizer are common infrastructure.}
\label{fig:method}
\end{figure*}
Let $\theta$ denote fixed model weights, $h_{\mathrm{in}}$ the phase incoming
harness, and $C=\{(x_i,x'_i,q_i,y_i,y'_i)\}$ valid answer-changing image pairs.
With correctness indicators $c_i$ and $c'_i$, define
\begin{equation}
 P_C(\theta,h)=|C|^{-1}\sum_i c_i c'_i. \label{eq:paired}
\end{equation}
The eligible candidate set is
\begin{equation}
 \mathcal H_{\mathrm{ok}}=\{h:P_C(\theta,h)\geq P_C(\theta,h_{\mathrm{in}})-\epsilon\}
 \cup\{h_{\mathrm{in}}\}. \label{eq:gate}
\end{equation}
We fix $\epsilon=0$ and select lexicographically:
\begin{equation}
 h_{\mathrm{next}}=\arg\max_{h\in\mathcal H_{\mathrm{ok}}}
       (A_H(\theta,h),-\mathrm{turns}_H(\theta,h)). \label{eq:select}
\end{equation}
Exact ties retain the original archive order. A subsequent native RSFT stage
uses successful trajectories under $h_{\mathrm{next}}$; its objective and success
filter are unchanged. These empirical constraints do not guarantee independent
test accuracy or non-regression after training.

The answer parser and verifier are shared across all conditions and cannot be
changed by candidates. Candidate code runs with file and network confinement.
Current weights, incoming harness, audit sample identities and decoding settings
are fixed within a selection phase.

To separate the gate from an ordinary accuracy check, let $M_C$ be single-image
accuracy on the same paired data and $q_0$ the fraction of pairs with both
answers incorrect. Expanding $(1-c_i)(1-c'_i)$ gives the exact identity
\[
 P_C=2M_C-1+q_0,\qquad \Delta P_C=2\Delta M_C+\Delta q_0.
\]
Thus equal single-image accuracy need not imply equal paired correctness.
If $q_0$ is constant across the fixed-weight candidate archive, the two
zero-tolerance gates are equivalent on that archive. Conversely, variation
in $q_0$ alone does not establish a different selected harness or a downstream
benefit. At fixed $M_C$, larger $P_C$ also entails more both-incorrect pairs;
the joint endpoint must therefore be interpreted alongside ordinary accuracy,
rather than as a guarantee that every type of error decreases.
Joint correctness does not identify why an answer fails. Perception, arithmetic,
instruction following and answer formatting can all affect this endpoint;
paired failures alone do not establish visual non-use or a latent mechanism.

\section{Experimental Protocol}
We use the public post-trained Qwen3.5-4B checkpoint~\cite{qwen35} and
seeds42,43,44. Each weight stage has four batches, eight
questions per batch and eight trajectories per question. Each search attempts
at most three candidates. Identical first stages and candidate measurements
are shared within a seed; subsequent branches are evaluated as dependent,
matched comparisons rather than independent first-stage trials.

PlotQA-EvidencePairs is a separately named dataset derived from PlotQA source
tables~\cite{plotqa}, with comparison, value-reading and difference questions.
All variants of a source table remain
in the same partition: W2048, H128, C64 pairs, V256 pairs, T1024 pairs, and an
independent recheck pool. C is optimization data. Numeric answers use a declared
tolerance; answer-changing intervals must not overlap. The external ChartQA~\cite{chartqa}
evaluation uses a preregistered subset of256 human and256 augmented questions;
it is not the full official benchmark score.

The common initial prompt and learning rate are selected using ordinary
development accuracy only. The final checkpoint is determined by the frozen
budget. Test data are excluded from all model and harness selection.

\section{Results}
\begin{table}[t]\centering\small
\begin{tabular}{lcc}\toprule
Condition & Single image & Both images\\\midrule
Weight only & --- & ---\\
Harness only & --- & ---\\
WHALE, one cycle & --- & ---\\
VETO, one cycle & --- & ---\\\bottomrule
\end{tabular}
\caption{Template only. No formal condition has a completed scientific result.
Final tables will include all seeds, mean, sample standard deviation and
source-cluster confidence intervals.}\label{tab:main}
\end{table}

The same-data marginal gate and counterfactual augmentation are mandatory
controls. If gates select the same harness, that observation provides no
evidence for a distinct benefit of VETO. We check the error decomposition on
the actual fixed-weight C archives; calibration measurements at different
weights cannot establish that the acceptance rules choose different harnesses.

\begin{table}[t]\centering\small
\begin{tabular}{lcc}\toprule
Selection / control & Paired score & Difference\\\midrule
Ordinary accuracy & --- & ---\\
Same-C marginal gate & --- & ---\\
VETO paired gate & --- & ---\\
50\% paired augmentation & --- & ---\\\bottomrule
\end{tabular}
\caption{Mandatory controls, pending real continuation experiments. Equal decisions
will be identified explicitly.}\label{tab:ablation}
\end{table}

\begin{table}[t]\centering\small
\begin{tabular}{lccc}\toprule
Condition & Human & Augmented & Cost\\\midrule
WHALE & --- & --- & ---\\
VETO & --- & --- & ---\\
Controls & --- & --- & ---\\\bottomrule
\end{tabular}
\caption{ChartQA fixed subset and total GPU/API costs. No external model score
is available in this scaffold.}\label{tab:external}
\end{table}

Stage trajectories and performance--cost plots will be generated from complete
branch records. Empty plots or simulated improvements are not experimental evidence.

\section{Limitations and Decision Criteria}
The study is limited to one model size, one domain and one complete alternation.
Source-derived renderings do not establish general visual-agent reliability.
The native checkpoint protocol omits Adam moments; all branches use the same
restoration behavior. We report full allocation costs including failed runs,
loading, candidate evaluation and audit overhead.

The prospective candidate criterion requires at least one percentage point
mean paired improvement, nonnegative differences across all three seeds,
positive differences in at least two seeds, ordinary accuracy loss no greater
than one point, and a positive lower bound of the paired source-cluster95\%
interval. Independent value beyond a marginal gate must also be demonstrated.
Unsupported outcomes will be reported as inconclusive rather than hidden.

\begin{thebibliography}{9}
\bibitem{whale} Haechan Kim et al. WHALE: A Simple Recipe for Joint Harness-Weight Optimization. arXiv:2609.00196,2026.

\bibitem{autodesign} Yaxin Luo et al. AutoDesign: Meta-Harness Optimization for Long-Horizon Agentic Design. arXiv:2608.13560,2026.

\bibitem{chartographer} Yifan Jiang et al. Chartographer: Counterfactual Chart Generation for Evaluating Vision-Language Models. arXiv:2605.27311,2026.

\bibitem{chartqa} Ahmed Masry et al. ChartQA: A Benchmark for Question Answering about Charts with Visual and Logical Reasoning. Findings of ACL,2022.

\bibitem{plotqa} Nitesh Methani, Pritha Ganguly, Mitesh M. Khapra, and Pratyush Kumar. PlotQA: Reasoning over Scientific Plots. WACV,2020.

\bibitem{qwen35} Qwen Team. Qwen3.5: Towards Native Multimodal Agents.2026. \href{https://qwen.ai/blog?id=qwen3.5}{qwen.ai/blog?id=qwen3.5}.

\end{thebibliography}
\end{document}
